\documentclass[11pt,letterpaper]{article}

% ── Packages ──────────────────────────────────────────────
\usepackage[margin=1in]{geometry}
\usepackage{enumitem}
\usepackage{hyperref}
\usepackage[dvipsnames]{xcolor}
\usepackage{microtype}
\usepackage{parskip}
\usepackage{titlesec}
\usepackage{booktabs}
\usepackage{fancyvrb}

% ── Formatting ────────────────────────────────────────────
\hypersetup{colorlinks=true, linkcolor=NavyBlue, citecolor=NavyBlue, urlcolor=NavyBlue}
\titleformat{\section}{\large\bfseries}{\thesection}{1em}{}
\titleformat{\subsection}{\normalsize\bfseries}{\thesubsection}{1em}{}
\setlength{\parskip}{0.6em}

\begin{document}

% ── Title ─────────────────────────────────────────────────
\begin{center}
{\LARGE\bfseries Cross-Patient Speech Decoding:\\[4pt]A Foundation Model Approach}\\[12pt]
{\normalsize Ben Tang \quad\textbar\quad Cogan Lab, Duke University \quad\textbar\quad April 2026}
\end{center}

\medskip\noindent\rule{\textwidth}{0.4pt}\medskip


%% ====================================================================
\section{The Problem}
%% ====================================================================

We have intra-operative \textmu ECoG recordings from ${\sim}$10 patients during a non-word repetition task (52 CVC/VCV tokens, 9 phonemes). Each patient has a thin-film electrode array on left sensorimotor cortex---128 or 256 channels at 1.3--1.7\,mm pitch. The goal: decode the 3-phoneme sequence the patient is saying from neural activity.

The fundamental challenge is that \textbf{every patient's array sits in a different place}. Arrays are offset 15--25\,mm from each other in standardized brain coordinates, with variable rotation. Channel 47 in Patient A records from a completely different patch of cortex than channel 47 in Patient B.

This means we can't simply pool data across patients---the spatial correspondence doesn't exist at the electrode level. Current approaches either train separate models per patient (missing shared structure) or use linear alignment methods like PCA+CCA (Spalding 2025, 31\% balanced accuracy)---a reasonable first pass, but one that doesn't exploit the spatial correspondence between arrays.

\textbf{What we want:} A single model trained on all patients that outperforms individual per-patient models, and that can be quickly calibrated to a new patient with minimal data.


%% ====================================================================
\section{The Insight}
%% ====================================================================

Think of each patient's electrode array as a \textbf{camera pointed at the brain from a slightly different angle}. Each camera sees a partial, noisy view of the same underlying scene---the neural dynamics of speech production. The scene is shared; the viewpoints differ.

In computer vision, multi-view reconstruction solves exactly this: given images from different cameras, recover the 3D scene. The recipe is:

\begin{enumerate}[leftmargin=2em]
    \item \textbf{Calibrate each camera} (intrinsics/extrinsics)
    \item \textbf{Define sample points} in 3D space where you want to reconstruct
    \item \textbf{Aggregate information} from whichever cameras can see each point
\end{enumerate}

We do the same thing with brains:

\begin{enumerate}[leftmargin=2em]
    \item \textbf{Per-patient calibration}---a small set of learned parameters (182 per patient) that absorb differences in signal gain/impedance and electrode registration error. This is the ``camera calibration.''

    \item \textbf{Virtual electrodes (VEs)}---16 fixed positions in standardized brain space, chosen from a neuroanatomical atlas (Brainnetome, derived from 40 healthy subjects). These positions correspond to known speech-relevant regions: tongue motor cortex, face motor cortex, Broca's area, auditory cortex, etc. These are the ``3D sample points.''

    \item \textbf{Cross-attention}---each virtual electrode ``looks at'' the real electrodes nearby and extracts a weighted summary. Electrodes close to a VE contribute more (distance bias), and the model learns \textit{what} to extract from them (shared attention weights). This is the ``multi-view aggregation.''
\end{enumerate}

The result: every patient's variable-count, variable-position electrode array gets mapped into the \textbf{same 16-dimensional functional space}. A tongue motor cortex VE always represents tongue motor cortex activity, regardless of which patient's electrodes contributed to it.


%% ====================================================================
\section{Why This Should Work}
%% ====================================================================

Three lines of evidence:

\subsection*{1.~Speech motor cortex has conserved organization across people}

The gross layout---lips, jaw, tongue, larynx arranged along the cortical surface---is the same in every human (Breshears 2015, $n{=}33$; Bouchard 2013; Roux 2020). The \textit{ordinal} arrangement is invariant; the \textit{absolute positions} vary by 5--15\,mm. Our architecture handles this: the atlas provides the average positions, and per-patient offsets (leashed to ${\leq}$15\,mm) adapt to individual anatomy.

\subsection*{2.~Multiple independent groups have converged on this design}

At least 5 groups have independently arrived at Perceiver-style cross-attention bottlenecks for multi-patient intracranial data:

\begin{itemize}[leftmargin=2em]
    \item \textbf{MIBRAIN} (Wu 2025)---sEEG, cross-patient consonant decoding via atlas region mapping
    \item \textbf{BarISTA} (Oganesian 2025, NeurIPS)---atlas-level encoding beats channel-level by 8--10 percentage points
    \item \textbf{Charmander} (Mahato 2025, NeurIPS WS)---Perceiver bottleneck for multi-patient iEEG
    \item \textbf{POYO} (Azabou 2023, NeurIPS)---motor decoding across 7 monkeys
    \item \textbf{Population Transformer} (Chau 2025, ICLR)---cross-patient sEEG with coordinate encoding
\end{itemize}

None has applied this to \textmu ECoG speech, and none combines atlas-grounded positions with per-patient adaptation---but the convergence validates the core mechanism.

\subsection*{3.~Scale without spatial mechanism fails}

The Brant family of models (Zhang 2023, NeurIPS) trained on up to 40,000 hours of intracranial data with up to 1 billion parameters---and fails at cross-patient speech decoding. NDT3 (Ye 2025) used 2,000 hours with 350M parameters and found cross-subject transfer fundamentally limited by sensor variability. These are the clearest negative results in the field: \textbf{raw scale is insufficient}. You need an explicit mechanism to handle electrode placement variation.


%% ====================================================================
\section{The Architecture}
%% ====================================================================

\begin{Verbatim}[frame=single, fontsize=\small]
Raw HGA signal (per electrode, 200 Hz)
  → Temporal binning (Conv1d, 50ms windows)
  → Per-patient gain/offset normalization (128 params/patient)
  → Add spatial identity (Fourier encoding of MNI coordinates)
  → VE cross-attention (16 atlas positions, distance-biased)
  → 2 blocks of [spatial self-attention + temporal self-attention]
  → Autoregressive decoder: predict phoneme 1, then 2, then 3
\end{Verbatim}

\textbf{${\sim}$175,000 shared parameters. 182 per-patient parameters.}

For context: BIT (the Utah array foundation model, Zhang 2026) uses ${\sim}$7 million parameters and 6 million per patient. Willett's clinical system uses ${\sim}$66,000 per session. We use 182.

The per-patient parameters do three things:

\begin{itemize}[leftmargin=2em]
    \item \textbf{Signal normalization} (128 params): absorbs impedance/gain differences. Initialized from each patient's channel statistics---starts as z-scoring, learns residual corrections.
    \item \textbf{Coordinate correction} (6 params): a small rigid translation + rotation of the electrode positions, absorbing systematic registration error (${\sim}$4--8\,mm from surgical reconstruction).
    \item \textbf{VE position offsets} (48 params): shifts each of the 16 atlas positions by up to 15\,mm to match the patient's individual somatotopy.
\end{itemize}

Everything else---the attention weights, the temporal processing, the decoder---is shared across all patients. The shared weights learn \textit{what speech activity looks like}; the per-patient weights learn \textit{where to find it in this patient}.


%% ====================================================================
\section{Training Plan}
%% ====================================================================

\subsection*{Phase 1 (immediate): Supervised cross-patient training on our data}

Train on $N{-}1$ patients, evaluate on the held-out patient. Establish whether the architecture produces a useful cross-patient model at all. Compare against the per-patient baseline (PER\,=\,0.734 on our best patient, S14).

This phase requires no new data---just implementation and compute on DCC.

\subsection*{Phase 2 (contingent on data access): Foundation model with external pretraining}

This is where the project becomes a foundation model paper. The idea:

\begin{enumerate}[leftmargin=2em]
    \item \textbf{Pretrain} on a large corpus of chronic ECoG speech recordings using self-supervised learning (predict masked neural activity). No phoneme labels needed---just raw neural data during speech.
    \item \textbf{Fine-tune} on our labeled CoganLab data (per-patient, 46--178 trials each).
\end{enumerate}

The architecture is \textbf{modality-agnostic}---it treats any intracranial recording as ``electrodes with positions and time series.'' This means we can pretrain on clinical ECoG (10\,mm electrode spacing) and fine-tune on our research \textmu ECoG (1.3\,mm spacing). The VE bottleneck abstracts away electrode density.

\textbf{This cross-density transfer is the unique story no other lab can tell alone.}


%% ====================================================================
\section{What We Need: External Data}
%% ====================================================================

The single highest-leverage action is obtaining external chronic ECoG speech data for pretraining. Two sources:

\textbf{Priority 1: Flinker/Chen lab (NYU).} 48 patients, chronic epilepsy monitoring ECoG, speech production tasks. Code is already public (Chen 2024, \textit{Nature Machine Intelligence}). Data available upon request with a standard sharing agreement. This is a straightforward PI-to-PI data request.

\textbf{Priority 2: Chang lab (UCSF).} ${\sim}$15--25 patients, high-density 256-channel ECoG, speech production. Papers state ``data available from corresponding author.'' This is better framed as a \textbf{collaboration}: three institutions (Duke, NYU, UCSF) contributing three electrode densities (1.3\,mm, 10\,mm, 4\,mm) to build the first cross-density intracranial foundation model. Each lab needs the others to tell the cross-density story.

Combined, this would give us ${\sim}$50--100 hours of diverse speech-related neural data for pretraining---100$\times$ more than our CoganLab data alone.

\textbf{Without external data}, we fall back to pretraining on our own smaller corpus (${\sim}$24 hours of CoganLab sEEG + \textmu ECoG). The architecture still works, but the paper becomes a methods paper instead of a foundation model paper.


%% ====================================================================
\section{The Paper}
%% ====================================================================

\textbf{Title direction:} \textit{Neural Field Perceiver: A Foundation Model for Intracranial Field Potentials}

\textbf{Core claim:} The first foundation model for intracranial high-gamma activity (HGA). BIT (Zhang 2026) built a foundation model for Utah array spikes; we do the same for field potentials. These are fundamentally different signals---HGA reflects postsynaptic potentials (inputs to cortex), spikes reflect action potentials (outputs). Different biophysics, different clinical modalities, no existing foundation model.

\textbf{Key results we expect to show:}

\begin{enumerate}[leftmargin=2em]
    \item Cross-patient v12 beats per-patient baselines (the architecture works)
    \item Performance scales with number of patients (the pooling works)
    \item Atlas-grounded VEs beat learned latents at small $N$ (the atlas prior matters)
    \item SSL pretraining shifts the calibration curve left (fewer trials needed per new patient)
    \item Cross-density transfer works (clinical ECoG pretraining helps research \textmu ECoG)
\end{enumerate}

\textbf{Worst case:} If nothing transfers, that's still a publishable negative result showing why intracranial field potential foundation models are hard. The systematic evaluation (8 experiment phases, ${\sim}$44 controlled ablations) would be the most thorough investigation of cross-patient intracranial transfer to date.


%% ====================================================================
\section{Current Status}
%% ====================================================================

\begin{itemize}[leftmargin=2em]
    \item Architecture fully designed (22-page technical design document)
    \item Per-patient baseline established (PER\,=\,0.734, S14)
    \item 55 prior cross-patient experiments characterizing the problem
    \item All electrode coordinates verified and mapped to MNI space (11 patients)
    \item 16 atlas positions selected and validated for patient reachability
    \item DCC compute ready (8$\times$ RTX 5000 Ada)
\end{itemize}

\textbf{Immediate blockers:}

\begin{enumerate}[leftmargin=2em]
    \item \textbf{External data request}---Greg needs to initiate PI-to-PI contact with Flinker (NYU) and/or Chang (UCSF).
    \item \textbf{Nonlinear MNI normalization}---current coordinate transform is linear-only (${\sim}$5--15\,mm residual error). If T1 MRIs exist, ANTs nonlinear warp could halve this. Need to check with Zac.
    \item \textbf{Implementation}---architecture is designed but not yet coded. Estimate ${\sim}$2 weeks to first training run on DCC.
\end{enumerate}


%% ====================================================================
\section*{Summary}
%% ====================================================================

Each patient's electrode array is a partial view of shared speech motor dynamics. We use anatomical atlas positions to define a common coordinate system, map each patient's electrodes into that space via cross-attention, and learn shared temporal representations that transfer across patients. 175K shared parameters, 182 per patient. The architecture is validated by convergent independent work from 5+ groups. What's missing is scale---and that requires data from other labs.


\end{document}
